% GENERATED FILE. Edit canonical lessons, then run npm run generate:books.
% canonical-source-hash: fnv1a64:16d1fca56615626e
% canonical-lessons: SA-C230-saphalata, SA-C230-viphalata, SA-C230-prayatnah, SA-C230-avasarah, SA-C230-anubhavah

\chapter{Success, Failure, Effort, Opportunity, Experience}
\label{ch:sa-success-failure-effort-opportunity-experience}
\hlchaptermodality{230}

\begin{chapteropening}
\textbf{By the end of this chapter:} \emph{I can say success, failure, an effort, an opportunity and experience.}

Complete the last lesson of chapter 230: I can say success, failure, an effort, an opportunity and experience.
\end{chapteropening}

\section[\texorpdfstring{\sk{सफलता} (saphalatā)}{सफलता (saphalatā)}]{\sk{सफलता} (saphalatā) — success}

\label{lesson:SA-C230-saphalata}



\begin{quote}
Before the new one: say the Sanskrit for authority, then the Sanskrit for a question paper.
\end{quote}

\subsection*{You'll want to know: \sk{सफलता}}
\textbf{\sk{सफलता}} — \emph{saphalatā} — \textquotedblleft{}success\textquotedblright{}.

\begin{grammarlens}[title={What you've built}]
One more word for this chapter.
\end{grammarlens}

\subsection*{Your turn}
\begin{itemize}
\raggedright
  \item \emph{Say it:} \emph{saphalatā}
  \item \emph{Say it:} \emph{saphalatā}, once more
  \item \emph{Say it:} say \emph{saphalatā}
  \item \emph{Recall:} say the Sanskrit for authority, then the Sanskrit for a question paper, then say \emph{saphalatā} again
\end{itemize}

\begin{tcolorbox}[breakable,colback=teal!4,colframe=teal!35!black,title={Before you move on}]
What does \sk{सफलता} mean? (\textquotedblleft{}Success\textquotedblright{}.) Say it once more.
\end{tcolorbox}
\section[\texorpdfstring{\sk{विफलता} (viphalatā)}{विफलता (viphalatā)}]{\sk{विफलता} (viphalatā) — failure}

\label{lesson:SA-C230-viphalata}



\begin{quote}
Before the new one: say the Sanskrit for a question paper, then the Sanskrit for success.
\end{quote}

\subsection*{You'll want to know: \sk{विफलता}}
\textbf{\sk{विफलता}} — \emph{viphalatā} — \textquotedblleft{}failure\textquotedblright{}.

\begin{grammarlens}[title={What you've built}]
One more word for this chapter.
\end{grammarlens}

\subsection*{Your turn}
\begin{itemize}
\raggedright
  \item \emph{Say it:} \emph{viphalatā}
  \item \emph{Say it:} \emph{viphalatā}, once more
  \item \emph{Say it:} say \emph{viphalatā}
  \item \emph{Recall:} say the Sanskrit for a question paper, then the Sanskrit for success, then say \emph{viphalatā} again
\end{itemize}

\begin{tcolorbox}[breakable,colback=teal!4,colframe=teal!35!black,title={Before you move on}]
What does \sk{विफलता} mean? (\textquotedblleft{}Failure\textquotedblright{}.) Say it once more.
\end{tcolorbox}
\section[\texorpdfstring{\sk{प्रयत्नः} (prayatnaḥ)}{प्रयत्नः (prayatnaḥ)}]{\sk{प्रयत्नः} (prayatnaḥ) — an effort}

\label{lesson:SA-C230-prayatnah}



\begin{quote}
Before the new one: say the Sanskrit for success, then the Sanskrit for failure.
\end{quote}

\subsection*{You'll want to know: \sk{प्रयत्नः}}
\textbf{\sk{प्रयत्नः}} — \emph{prayatnaḥ} — \textquotedblleft{}an effort\textquotedblright{}.

\begin{grammarlens}[title={What you've built}]
One more word for this chapter.
\end{grammarlens}

\subsection*{Your turn}
\begin{itemize}
\raggedright
  \item \emph{Say it:} \emph{prayatnaḥ}
  \item \emph{Say it:} \emph{prayatnaḥ}, once more
  \item \emph{Say it:} say \emph{prayatnaḥ}
  \item \emph{Recall:} say the Sanskrit for success, then the Sanskrit for failure, then say \emph{prayatnaḥ} again
\end{itemize}

\begin{tcolorbox}[breakable,colback=teal!4,colframe=teal!35!black,title={Before you move on}]
What does \sk{प्रयत्नः} mean? (\textquotedblleft{}An effort\textquotedblright{}.) Say it once more.
\end{tcolorbox}
\section[\texorpdfstring{\sk{अवसरः} (avasaraḥ)}{अवसरः (avasaraḥ)}]{\sk{अवसरः} (avasaraḥ) — an opportunity}

\label{lesson:SA-C230-avasarah}



\begin{quote}
Before the new one: say the Sanskrit for failure, then the Sanskrit for an effort.
\end{quote}

\subsection*{You'll want to know: \sk{अवसरः}}
\textbf{\sk{अवसरः}} — \emph{avasaraḥ} — \textquotedblleft{}an opportunity\textquotedblright{}.

\begin{grammarlens}[title={What you've built}]
One more word for this chapter.
\end{grammarlens}

\subsection*{Your turn}
\begin{itemize}
\raggedright
  \item \emph{Say it:} \emph{avasaraḥ}
  \item \emph{Say it:} \emph{avasaraḥ}, once more
  \item \emph{Say it:} say \emph{avasaraḥ}
  \item \emph{Recall:} say the Sanskrit for failure, then the Sanskrit for an effort, then say \emph{avasaraḥ} again
\end{itemize}

\begin{tcolorbox}[breakable,colback=teal!4,colframe=teal!35!black,title={Before you move on}]
What does \sk{अवसरः} mean? (\textquotedblleft{}An opportunity\textquotedblright{}.) Say it once more.
\end{tcolorbox}
\section[\texorpdfstring{\sk{अनुभवः} (anubhavaḥ)}{अनुभवः (anubhavaḥ)}]{\sk{अनुभवः} (anubhavaḥ) — experience}

\label{lesson:SA-C230-anubhavah}



\begin{quote}
Before the new one: say the Sanskrit for an effort, then the Sanskrit for an opportunity.
\end{quote}

\subsection*{You'll want to know: \sk{अनुभवः}}
\textbf{\sk{अनुभवः}} — \emph{anubhavaḥ} — \textquotedblleft{}experience\textquotedblright{}.

\begin{grammarlens}[title={What you've built}]
One more word for this chapter.
\end{grammarlens}

\subsection*{Your turn}
\begin{itemize}
\raggedright
  \item \emph{Say it:} \emph{anubhavaḥ}
  \item \emph{Say it:} \emph{anubhavaḥ}, once more
  \item \emph{Say it:} say \emph{anubhavaḥ}
  \item \emph{Recall:} say the Sanskrit for an effort, then the Sanskrit for an opportunity, then say \emph{anubhavaḥ} again
\end{itemize}

\begin{tcolorbox}[breakable,colback=teal!4,colframe=teal!35!black,title={Before you move on}]
What does \sk{अनुभवः} mean? (\textquotedblleft{}Experience\textquotedblright{}.) Say it once more.
\end{tcolorbox}
